\documentclass[twocolumn]{article}
\usepackage[utf8]{inputenc}
\usepackage{amsmath}
\usepackage{amssymb}
\usepackage{graphicx}

\title{Supplementary Methods: Phenotype-First Evolution and Extended Agency}
\author{Anonymous}
\date{}

\begin{document}

\maketitle

% ════════════════════════════════════════════════════════════════════
% SUPPLEMENTARY METHODS
% ════════════════════════════════════════════════════════════════════

\section*{Detailed Algorithm Descriptions}

\subsection*{Gene-Centric GA (Baseline)}

The Gene-Centric algorithm represents the minimal evolutionary baseline: pure mutation without learning.

\textbf{Representation}: Each organism has a binary genome of length 100 bits, where each bit represents whether a particular cell (on the 10$\times$10 grid) is occupied.

\textbf{Fitness evaluation}: The organism's phenotype is its genome. Fitness is computed as described in Methods: mean recall across all patterns in the task.

\textbf{Selection}: Tournament selection with tournament size 3. The best individual in each tournament is selected for breeding.

\textbf{Reproduction}: Selected individuals are copied. Each copy undergoes mutation: each bit is flipped with probability 2\% (uniform random mutation).

\textbf{Population dynamics}: Population size 50, 150 generations. No learning, no inheritance of learned information.

\textbf{Convergence criterion}: None. Algorithm runs for full 150 generations.

\subsection*{Baldwin Effect (Learning Without Inheritance)}

The Baldwin Effect tests whether selection-based learning suffices for evolution.

\textbf{Learning phase}: Each organism in each generation is given 20 learning trials. In each trial:
\begin{enumerate}
\item A random subset of cells is activated (set to 1) with probability 30\% per cell
\item Fitness is computed for this modified phenotype
\item The organism keeps the change if fitness improves; otherwise, reverts
\item This continues for 20 trials, allowing gradient-free hill-climbing
\end{enumerate}

After 20 trials, the organism has a learned phenotype that may be better than its genome.

\textbf{Fitness evaluation}: The organism's fitness for selection is its learned phenotype fitness (after 20 learning trials).

\textbf{Selection}: Tournament selection (size 3) based on learned fitness.

\textbf{Reproduction}: Selected individuals are copied. Offspring genomes are identical to parents' genomes (no learning is inherited). Each offspring undergoes mutation (2\% bit flip rate). Offspring will then learn independently.

\textbf{Key point}: Learned information is NOT passed to offspring. Each generation's organisms must learn from scratch. Selection favors organisms whose genomes permit better learning, but learned patterns themselves are not inherited.

\subsection*{Lamarckian GA (Slow Genetic Assimilation)}

The Lamarckian algorithm tests slow write-back via genetic assimilation.

\textbf{Learning phase}: Identical to Baldwin Effect. Each organism learns via 20 trials.

\textbf{Fitness evaluation}: Learned phenotype fitness.

\textbf{Selection}: Tournament selection (size 3).

\textbf{Genetic assimilation}: After selection, before reproduction, a fraction of learned information is encoded into the selected organism's genome.

Specifically, for each cell in the learned phenotype:
\begin{itemize}
\item If the cell is occupied in learned phenotype (value 1): set genome bit to 1 with probability 15\%
\item If the cell is empty in learned phenotype (value 0): set genome bit to 0 with probability 15\%
\end{itemize}

This 15\% encoding rate means that after 1 generation, 15\% of learned information enters the genome. After 2 generations (if the same organism is selected), up to 15\% $\times$ 2 = 30\% is encoded. Full encoding takes approximately 6--7 generations for a single trait, consistent with genetic assimilation timescales of 50--100 generations for whole populations.

\textbf{Reproduction}: Selected organisms are copied with their updated genomes (now partially encoding learned information). Offspring undergo mutation (2\% per bit).

\textbf{Key point}: Learned information slowly encodes into genes via selection pressure. This is the classic genetic assimilation mechanism documented by Waddington and formalized by Hinton and Nowlan.

\subsection*{Phenotype-First GA (Epigenomic Inheritance)}

The Phenotype-First algorithm tests direct, fast write-back via epigenomic inheritance.

\textbf{Organism structure}: Each organism has:
\begin{itemize}
\item A genome (binary 100-bit string, initial random)
\item An epigenome (binary 100-bit string, initially empty)
\end{itemize}

The phenotype is the union of genome and epigenome: phenotype[i] = max(genome[i], epigenome[i]). This represents open chromatin (epigenome) overlaid on genetic instructions.

\textbf{Learning phase}: Each organism learns via 20 trials on a temporary phenotype copy. Learning modifies the epigenome directly (not the genome). After learning, the organism has:
\begin{itemize}
\item Genome (unchanged from birth)
\item Epigenome (updated from learning)
\end{itemize}

\textbf{Fitness evaluation}: Fitness is computed on the learned phenotype (genome $\cup$ epigenome).

\textbf{Selection}: Tournament selection (size 3) based on learned fitness.

\textbf{Epigenomic inheritance}: Selected organisms are reproduced. Offspring receive:
\begin{itemize}
\item A mutated copy of parent's genome (2\% per bit)
\item A complete copy of parent's epigenome (direct copy, 100\% fidelity, no mutation)
\end{itemize}

This is the key difference: learned patterns in the epigenome are passed directly to offspring without waiting for genetic assimilation. Offspring inherit the parent's accumulated learning.

\textbf{Key point}: Direct, fast write-back (1 generation vs. 50--100). Offspring benefit from parent's learned patterns immediately.

% ════════════════════════════════════════════════════════════════════
% STATISTICAL METHODS
% ════════════════════════════════════════════════════════════════════

\section*{Statistical Analysis}

\subsection*{Comparison Between Methods}

For each task (simple or compositional), we compared fitness distributions across methods using:

\textbf{ANOVA}: One-way analysis of variance testing the null hypothesis that all methods have equal mean fitness. We report F-statistic and p-value.

\textbf{Welch's t-test}: Pairwise comparison between Phenotype-First and Baldwin Effect, not assuming equal variances. Reports t-statistic and p-value.

\textbf{Cohen's d}: Effect size between Phenotype-First and Baldwin. Interpreted as:
\begin{itemize}
\item $d < 0.2$: negligible effect
\item $0.2 \leq d < 0.5$: small effect
\item $0.5 \leq d < 0.8$: medium effect
\item $d \geq 0.8$: large effect
\end{itemize}

\textbf{Standard deviation}: Reported as $\pm$ notation. For Phenotype-First on T+Plus, we report $s.d. = 0.0$, meaning all 30 runs achieved identical fitness values. This indicates deterministic, non-random learning.

\subsection*{Success Rate}

We define ``success rate'' as the fraction of runs achieving fitness $> 0.8$ (80\% of maximum). This captures high-fidelity learning of patterns.

We also report ``both-present rate'' for compositional tasks: the fraction of runs where both patterns are present (pattern score $\geq 0.5$ for both components).

% ════════════════════════════════════════════════════════════════════
% CONVERGENCE ANALYSIS
% ════════════════════════════════════════════════════════════════════

\section*{Convergence Analysis}

For simple pattern learning, we measured convergence speed as the generation at which fitness first exceeded 0.8. Results:

\begin{itemize}
\item Baldwin Effect: $3.3 \pm 2.1$ generations
\item Phenotype-First: $2.7 \pm 1.9$ generations
\item Difference: not significant ($p = 0.559$)
\end{itemize}

For simple learning, both methods converge at similar speeds. The difference between 3.3 and 2.7 generations is within statistical noise.

For compositional learning, convergence is much slower and depends on the mechanism. Phenotype-First shows steady improvement due to accumulated inheritance, while Baldwin shows stochastic wandering without consistent accumulation.

% ════════════════════════════════════════════════════════════════════
% ROBUSTNESS CHECKS
% ════════════════════════════════════════════════════════════════════

\section*{Robustness and Sensitivity Analysis}

\subsection*{Mutation Rate}

We tested mutation rates from 0.5\% to 5\%. Results were robust across this range, with 2\% performing well for all algorithms. Higher mutation rates (5\%) slightly slowed convergence but did not qualitatively change results.

\subsection*{Population Size}

We tested population sizes 25, 50, and 100. Population size 50 was chosen as a balance between computational cost and convergence. Larger populations showed faster convergence but similar final fitness.

\subsection*{Learning Trial Count}

We tested learning trial counts from 5 to 50 trials per organism per generation. With 20 trials, learning was sufficient for good performance without excessive computation. More trials slightly improved performance on compositional tasks but showed diminishing returns.

\subsection*{Epigenomic Fidelity}

For Phenotype-First, we tested whether perfect fidelity (100\%) in epigenomic inheritance was necessary. We tested perfect copy (100\%) versus noisy copy (95\%, 90\%, 85\% fidelity). Results showed that even with 90\% fidelity, Phenotype-First outperformed Baldwin on compositional tasks, but perfect fidelity showed highest reproducibility.

% ════════════════════════════════════════════════════════════════════
% BIOLOGICAL PLAUSIBILITY
% ════════════════════════════════════════════════════════════════════

\section*{Biological Plausibility and Limitations}

Our model is a simplified abstraction. Real organisms have vastly more complex learning and genetic systems. Specifically:

\textbf{Direct copy assumption}: We assume epigenomic inheritance is a perfect copy (deepcopy in code). Real epigenomic inheritance involves DNA methylation, histone modifications, and other mechanisms that may degrade over generations. However, evidence for multi-generational epigenetic inheritance (particularly through stress-response pathways) supports the basic principle.

\textbf{Learning mechanism}: We model learning as a gradient-free random walk with fitness feedback. Real organisms use diverse learning mechanisms (neural networks, behavioral trial-and-error, metabolic adaptation). Our abstract model captures the essence: organisms can improve their phenotype through environmental interaction.

\textbf{Substrate}: We use a simple 10$\times$10 grid and geometric shapes. Real evolutionary problems involve complex ecological niches, developmental constraints, and trade-offs. Compositional tasks approximate real problems requiring integration of multiple learned skills (e.g., predator avoidance AND foraging).

\textbf{Trade-offs not modeled}: Real organisms face trade-offs between learning investment and reproduction, between fidelity and plasticity of inherited information. Our simplified model avoids these to isolate the write-back mechanism.

\end{document}
